% !TEX root = ../main.tex
\section{Introducción}

En la actualidad, las organizaciones manejan grandes cantidades de
información que deben mantenerse organizadas para facilitar su consulta,
actualización, almacenamiento y recuperación. En instituciones como las
universidades, esta necesidad adquiere especial importancia debido al
volumen de datos personales y académicos que deben administrarse de forma
constante.

Antes de desarrollar un sistema informático es necesario comprender el
problema que se desea resolver, identificar los datos que serán
administrados y definir estructuras adecuadas para representarlos. De
igual forma, se debe considerar la eficiencia de las operaciones que se
realizarán sobre dicha información, debido a que distintas soluciones
pueden requerir diferentes cantidades de tiempo y recursos.

La presente investigación analiza el caso de la Universidad XUNA, una
institución con aproximadamente 25\,000 estudiantes que busca reemplazar
el uso de formularios impresos, hojas electrónicas y programas
independientes por un Sistema de Gestión Académica centralizado. El
sistema deberá permitir administrar información personal y académica de
los estudiantes, además de realizar operaciones como registrar, buscar,
modificar, consultar y generar reportes.

A partir de este caso se estudiarán conceptos fundamentales relacionados
con las estructuras de datos, entre ellos el dato, los Tipos de Datos
Abstractos, los registros, los arreglos, las funciones, el costo
computacional y los mecanismos de almacenamiento mediante archivos
secuenciales e indexados. El objetivo es establecer una propuesta de
diseño organizada y técnicamente justificada que sirva como base para el
posterior desarrollo del sistema.